\section{Putting Everything Together}
\label{sec:together}

To construct an end-to-end bitwidth-aware SNARK for circuit
satisfiability (Circ-SAT), we combine the building blocks
of Section~\ref{sec:blocks} to prove the PLONK constraints
of Definition~\ref{def:plonk}.
The same approach can be adapted to other constraint systems,
such as R1CS.
In addition to witness values, the constraint system encodes
wire positions to certify the wiring of the circuit.
Bitwidth awareness must therefore account for both the
\(\lceil\log_2 q\rceil\)-bit witness entries and the
\(\lceil\log_2 N\rceil\)-bit position labels.
We seek prover time and memory usage that scale linearly
with this \emph{effective bitwidth} \(\lceil\log_2 q\rceil + \lceil\log_2 N\rceil\).
For large circuits, processing position labels can introduce
an additional \(\log N\) factor in the corresponding costs.
When packing multiple executions of the same circuit,
the shared wiring allows its encoding and the associated
proof work to be amortized across executions.

\subsection{Proof for the Copy Constraint}
\label{sec:plonk-copy}

We reduce the PLONK copy constraint
\(\sigma(\bm w)=\bm w\) in Eq.~\eqref{eq:plonk2} to a
permutation relation over an intermediate field.
Let \(\bm w\in\Field^L\), where
\(L=3N+n_{\mathsf{io}}\), and let \(q=|\Field|\).
The reduction encodes each witness value together with its
position.

\smallskip\noindent\textbf{Joint encoding of values and positions.}
Choose \(\underline q=2^s\leq q\), with
\(s=\Theta(\log_2q)\), and elements
\(\delta_1,\ldots,\delta_s\in\Field\) such that the map
\(\sum_{t=1}^s2^{t-1}b_t\mapsto\sum_{t=1}^s\delta_tb_t\)
is injective on \(\{0,1\}^s\).
For prime \(q\), take \(s=\lfloor\log_2q\rfloor\) and
\(\delta_t=2^{t-1}\).
Appendix~\ref{appx:mleiota} gives a choice for every finite field.
We reserve \(\lceil\log_{\underline q} L\rceil\) base-field
coordinates for the position and one for the witness value.
Choose \(d\) as the least power of two at least
\(\lceil\log_{\underline q}L\rceil+1\), and choose a
power-of-two extension degree \(\kappa\) with \(2d\leq\kappa\).
This gives the tower:
\[
    \Field\subset\mathbb K\subset\ExtensionField,
    \quad
    [\mathbb K:\Field]=d,
    \quad
    d\mid\kappa,
    \quad
    [\ExtensionField:\mathbb K]=\kappa_{\mathbb K}
       =\kappa/d.
\]
We fix an \(\Field\)-basis
\(\bm\theta=(\theta_1,\ldots,\theta_d)^\top\in\mathbb K^d\)
of \(\mathbb K\).

Then, we can represent \(\sigma\) by
\(\bm\sigma=(\sigma(1),\ldots,\sigma(L))^\top\), and let
\(\bm i=(1,\ldots,L)^\top\).
Write their zero-based positions in radix \(\underline{q}\):
\[
\begin{aligned}
    \bm\sigma-\bm1_L
       &=\sum_{j=1}^{\lceil\log_{\underline{q}} L\rceil}\underline q^{j-1}\bm\sigma_{(j)},
       \quad
    \bm i-\bm1_L
       =\sum_{j=1}^{\lceil\log_{\underline{q}} L\rceil}\underline q^{j-1}\bm i_{(j)},
    \quad
    \bm\sigma_{(j)},\bm i_{(j)}
       \in [0,\underline{q}-1]^L.
\end{aligned}
\]
These are integer decompositions. In field expressions, we
interpret each digit through the injection above.
Consequently,
we encode the permutation images and identity positions as:
\[
    \bm\varsigma
       \coloneqq\sum_{j=1}^{\lceil\log_{\underline q} L\rceil}\theta_j\bm\sigma_{(j)} \in \mathbb{K}^L,
    \qquad
    \bm\iota
       \coloneqq\sum_{j=1}^{\lceil\log_{\underline q} L\rceil}\theta_j\bm i_{(j)}
       \in\mathbb K^L.
\]
The labels \(\iota_1,\ldots,\iota_L\) are distinct, and
\(\varsigma_i=\iota_{\sigma(i)}\).
The remaining basis coordinate binds each label to a witness
value. We define:
\[
    \bm a\coloneqq\bm\varsigma+\theta_d \bm w  \in \mathbb{K}^L,
    \qquad
    \bm b\coloneqq\bm\iota+\theta_d\bm w
       \in\mathbb K^L.
\]

\begin{lemma}[Copy constraints from a permutation relation]
\label{eq:permutation}
For the encodings above and every \(\bm w\in\Field^L\):
\[
    \sigma(\bm w)=\bm w
    \quad\Longleftrightarrow\quad
    \bm a\stackrel{\mathsf{perm}}{=}\bm b.
\]
\end{lemma}

The basis coordinates make the joint encoding injective.
Consequently, the label in \(a_i\) can match only the label in
\(b_{\sigma(i)}\), and equality of these entries enforces
\(w_i=w_{\sigma(i)}\).
Appendix~\ref{appx:proof-copy-constraint} gives the proof.

\smallskip\noindent\textbf{Applying the packed permutation proof.}
With the witness commitments fixed, both operands are determined
by \(\bm w\) and the public permutation.
We treat \(\mathbb K\) as the base field and apply
Algorithm~\ref{alg:permutation} over \(\ExtensionField\)
with packing width \(\kappa_{\mathbb K}\), chosen so that
\(|\mathbb K|\geq2\kappa_{\mathbb K}-1\).
All required interpolation points can therefore be chosen
in \(\mathbb K\).
The permutation proof leaves weighted evaluations with arbitrary
column weights and tensor row weights on \(\bm a\) and \(\bm b\).
By linearity, these reduce to evaluations of
\(\bm\varsigma,\bm\iota\in\mathbb K^L\) and
\(\bm w\in\Field^L\).
The verifier evaluates \(\bm\iota\) directly using
\(O(\kappa_{\mathbb K}+\log L)\) extension-field operations,
as shown in Appendix~\ref{appx:proof-position-encoding}; ordinary
MLE queries need only \(O(\log L)\) operations.
The remaining claims are verified using the packed PCS.

\subsection{Packing for a Single Circ-SAT Instance}
\label{sec:plonk-packing}

With the copy constraints handled by the permutation proof,
it remains to prove the PLONK gate constraints in Eq.~\eqref{eq:plonk1}.
Let \(N=m\kappa\), with \(m,\kappa\) powers of two, and
reuse the polynomial basis of Section~\ref{sec:pcs}, with
\(\ext\beta_1=1\), for the wire commitments.
The packed Hadamard proof of Section~\ref{sec:hadamard}
handles the length-\(N\) gate equation.
For the copy proof, pad both operands with zeros to length
\(4N\) and sample its fingerprint challenge from
\(\ExtensionField\setminus\{0\}\).
The common padding factors then cancel.

We reduce the gate constraints to weighted Hadamard claims by
evaluating the residual's multilinear extension. The verifier samples
\(\ext{\bm x}\sample\ExtensionField^{\log_2N}\) and sets
\(\ext{\bm\xi}=\tensorspan{\ext{\bm x}}\).
It checks one random projection of the gate residual:
\[
    \left\langle \ext{\bm\xi},
    \bm q_L\circ\bm w_L+\bm q_R\circ\bm w_R
    +\bm q_M\circ(\bm w_L\circ\bm w_R)
    +\bm q_O\circ\bm w_O+\bm q_C\right\rangle=0.
\]
Define the virtual vector \(\bm w_M=\bm w_L\circ\bm w_R\).
The prover sends
\(\ext c_L,\ext c_R,\ext c_M,\ext c_O,\ext c_C\);
the verifier checks their sum and the following claims:
\begin{align}
    &\ext c_L+\ext c_R+\ext c_M+\ext c_O+\ext c_C=0,
    \nonumber\\
    &\ext c_L=\langle\ext{\bm\xi},\bm q_L\circ\bm w_L\rangle,
    \quad
    \ext c_R=\langle\ext{\bm\xi},\bm q_R\circ\bm w_R\rangle,
    \quad
    \ext c_O=\langle\ext{\bm\xi},\bm q_O\circ\bm w_O\rangle,
    \label{eq:plonkWeightedHadamard1}\\
    &\ext c_M=\langle\ext{\bm\xi},\bm q_M\circ\bm w_M\rangle,
    \quad
    \ext c_C=\langle\ext{\bm\xi},\bm q_C\rangle.
    \label{eq:plonkWeightedHadamard2}
\end{align}
Each of the three products in
Eq.~\eqref{eq:plonkWeightedHadamard1} has the form of
Eq.~\eqref{eq:hadamard-claim}. Algorithm~\ref{alg:hadamard}
reduces each to a selector evaluation and an evaluation of a committed
wire vector.

For \(\ext c_M\), the same reduction leaves an evaluation of
\(\bm q_M\) and an evaluation of \(\bm w_M\). The latter is
itself a weighted Hadamard claim on \(\bm w_L,\bm w_R\), which
we prove with a second invocation of Algorithm~\ref{alg:hadamard}.

We batch the terminal claims into a witness group and a public
group, using the packed opening reduction of Section~\ref{sec:pcs}.
The prover commits to the concatenation of the three packed wire
vectors before any challenge; the selector and position-table
commitment is preprocessed, with a total packed length of \(m_{\mathsf{pub}}=
(5+3\lceil\log_{\underline q}(3N+n_{\mathsf{io}})\rceil)m\).

\begin{theorem}[Single-instance packing]
\label{theorem:single-instance-packing}
Assume the parameters above, \(m\geq\max\{\kappa,q,n_{\mathsf{io}}\}\),
and the PCS assumptions of Proposition~\ref{proposition:packed-pcs}.
The single-instance protocol is complete, with soundness error
at most:
\begin{equation}
    \epsilon_\mathsf{single}
       =O\left(\frac{N}{|\ExtensionField|}\right) + \epsilon_{\mathsf{PCS}},
\label{eq:together-soundness}
\end{equation}
where \(\epsilon_{\mathsf{PCS}}\) bounds the joint binding and
opening error of the underlying PCS calls.
With circuit preprocessing, the efficiency bounds are:
\begin{equation}
\begin{aligned}
    \text{Prover time}\quad
       &O\left(N\kappa\log_q N\log_2\kappa\right)\,\bb
          +O(\kappa^4)
          +\mathcal{TP}_{\mathsf{pcs}}(3m+m_{\mathsf{pub}}),\\
    \text{Prover memory}\quad
       &O(N\log_q N)\,\Field+O(\kappa^4+n_{\mathsf{io}})
          +\mathcal{SP}_{\mathsf{pcs}}(3m+m_{\mathsf{pub}}),\\
    \text{Verifier time}\quad
       &O\left(
          \kappa\log_q N\log_2N+(\log_2N)^2
          +n_{\mathsf{io}}
          \right)\,\ee
          +\mathcal{TV}_{\mathsf{pcs}}(3m+m_{\mathsf{pub}}),\\
    \text{Proof size}\quad
       &O\left(\kappa\log_q N
          +\kappa\log_2\kappa
          +(\log_2N)^2\right)\,\ExtensionField
          +\mathcal{S}_{\mathsf{pcs}}(3m+m_{\mathsf{pub}}),
\end{aligned}
\label{eq:together-single-efficiency}
\end{equation}
where \(\mathcal{TP}_{\mathsf{pcs}}\),
\(\mathcal{SP}_{\mathsf{pcs}}\),
\(\mathcal{TV}_{\mathsf{pcs}}\), and
\(\mathcal{S}_{\mathsf{pcs}}\) denote the PCS
prover time, prover memory, verifier time, and proof size for
total logical length \(3m+m_\mathsf{pub}\), respectively,
including domain padding. The unlabeled remainders count
base-field operations and base-field elements.
\end{theorem}

The proof is deferred to Appendix~\ref{appx:proof-single-instance}.
Algorithm~\ref{alg:single-instance-packing} gives the protocol.

\subsection{Packing Multiple Circ-SAT Instances of the Same Circuit}
\label{sec:multi-instance-packing}

The $\log_q N$ term in
Theorem~\ref{theorem:single-instance-packing} introduced by the
position-encoding overhead can be amortized by packing executions of the same circuit.
Since all instances share the same selectors and wiring
permutation, the corresponding linear operations commute
with packing.
This yields amortized per-instance prover time
\(O(N\kappa)\,\bb\) and memory \(O(N)\,\Field\), excluding
the underlying PCS.

Consider \(\kappa\) executions of the same \(N\)-gate circuit,
where \(N,\kappa\) are powers of two,
\(\kappa\geq2\), and \(N\geq\max\{\kappa,n_{\mathsf{io}}\}\).
Let \(L=3N+n_{\mathsf{io}}\).
For each \(j\in[\kappa]\), let
\(\bm w^{(j)}\in\Field^{3N+n_{\mathsf{io}}}\)
be the witness vector from Definition~\ref{def:plonk}.
Using the polynomial basis \(\ext{\bm\beta}\) of
Section~\ref{sec:pcs}, with \(\ext\beta_1=1\), define:
\[
    \bm W
       =\concat{\bm w^{(1)},\ldots,\bm w^{(\kappa)}}
       \in\Field^{L\times\kappa},
    \qquad
    \ext{\bm w}=\bm W\ext{\bm\beta}\in\ExtensionField^L.
\]
For \(\ell\in\{L,R,O\}\), arrange the corresponding wire
vectors as \(\bm W_\ell\in\Field^{N\times\kappa}\).
Their packings and the virtual product vector are:
\begin{equation}
    \ext{\bm w}_\ell=\bm W_\ell\ext{\bm\beta}
       \in\ExtensionField^N,
    \qquad
    \ext{\bm w}_M=(\bm W_L\circ\bm W_R)\ext{\bm\beta}.
\label{eq:together-multiple-wires}
\end{equation}
The prover commits to the concatenation of
\(\ext{\bm w}_L,\ext{\bm w}_R,\ext{\bm w}_O\), of logical
length \(3N\), before any challenge.
The product vector remains uncommitted.

\begin{lemma}[Packed PLONK constraints]
\label{lemma:packedPlonk}
The \(\kappa\) instances satisfy the PLONK constraints if and only if:
\begin{align}
    \bm q_L\circ\ext{\bm w}_L
    +\bm q_R\circ\ext{\bm w}_R
    +\bm q_M\circ\ext{\bm w}_M
    +\bm q_O\circ\ext{\bm w}_O
    +\left(\sum_{j=1}^{\kappa}\ext\beta_j\right)\bm q_C
       &=\bm0_N,
       \label{eq:together-multiple-gates}\\
    \sigma(\ext{\bm w})&=\ext{\bm w}.
       \label{eq:packed-plonk2}
\end{align}
\end{lemma}

The shared selectors and permutation are \(\Field\)-linear,
and the basis packing is injective.
Thus each packed equality holds if and only if it holds in
every column of \(\bm W\).

\smallskip\noindent\textbf{Gate constraints.}
The verifier samples
\(\ext{\bm x}\sample\ExtensionField^{\log_2N}\)
and uses sumcheck to verify that the
\(\tensorspan{\ext{\bm x}}\)-weighted residual sum in
Eq.~\eqref{eq:together-multiple-gates} is zero.
The terminal check requires five selector MLE claims
and four packed-wire MLE evaluations at a common point
\(\ext{\bm y} \in \ExtensionField^{\log_2 N}\).
These claims are authenticated by the batched openings below.

To authenticate the packed product MLE value \(\ext v_M\),
the prover proves:
\begin{equation}
    \ext v_M
    =
    \tensorspan{\ext{\bm y}}^\top
    (\bm W_L\circ\bm W_R)\ext{\bm\beta}.
\label{eq:together-multiple-product-opening}
\end{equation}
The packed Hadamard-product proof in Algorithm~\ref{alg:hadamard}
reduces this to two weighted evaluation claims on
\(\bm W_L\) and \(\bm W_R\).
The column-projection step of Section~\ref{sec:pcs} converts
them to inner-product claims on the committed wire packings,
which are batched with the ordinary wire queries below.

\smallskip\noindent\textbf{Copy constraints.}
We apply the position encoding of Section~\ref{sec:plonk-copy}
in the basis \(\ext{\bm\beta}\) to obtain
\(\bm\iota,\bm\varsigma\in\ExtensionField^L\).
Appendix~\ref{appx:mleiota} evaluates the identity-label MLE
in \(O(\log L)\) extension-field operations.
The labels \(\iota_i\) are distinct and
\(\varsigma_i=\iota_{\sigma(i)}\).
For a fresh challenge \(\ext y\sample\ExtensionField\),
the copy constraint in Eq.~\eqref{eq:packed-plonk2} reduces
to a permutation check over $\ExtensionField$:
\[
    \bm\iota+\ext y\,\ext{\bm w}
       \stackrel{\mathsf{perm}}{=}
    \bm\varsigma+\ext y\,\ext{\bm w} \in \ExtensionField^L.
\]
The verifier samples an independent
\(\ext z\sample\ExtensionField\), and the prover uses two
extension-field product trees to establish:
\[
    \prod_{i=1}^{L}(\ext z+\iota_i+\ext y\,\ext w_i)
       =
    \prod_{i=1}^{L}(\ext z+\varsigma_i+\ext y\,\ext w_i).
\]
Each tree is padded with ones to length \(4N\).
The two reductions leave MLE claims on \(\bm\varsigma\),
\(\bm\iota\), and \(\ext{\bm w}\).
The verifier evaluates \(\bm\iota\) directly and computes the
public-I/O and padding contributions without commitments.
All remaining claims are batched as in
Section~\ref{sec:plonk-packing}, separately for the witness
and public vectors, with one underlying PCS opening per group.
When \(5+3\lceil\log_{\underline q}L\rceil\leq\kappa\),
the five selectors and the position-digit columns for the three
wire segments fit into one \(N\times\kappa\) base-field matrix.
Packing its rows gives one preprocessed public commitment of
length \(N\), for a total logical PCS length of \(4N\).

\begin{theorem}[Multiple-instance packing]
\label{theorem:multiple-instance-packing}
Assume the parameters above, \(q\geq2\kappa-1\),
\(5+3\lceil\log_{\underline q}L\rceil\leq\kappa\), and the PCS
assumptions of Proposition~\ref{proposition:packed-pcs}.
The protocol for \(\kappa\) instances of the same circuit
is complete, with soundness error at most:
\begin{equation}
    \epsilon_{\mathsf{multi}}
       =O\left(\frac{N}{|\ExtensionField|}\right)
          +\epsilon_{\mathsf{PCS}},
\label{eq:together-multiple-soundness}
\end{equation}
where \(\epsilon_{\mathsf{PCS}}\) bounds the joint binding and
opening error of the underlying PCS calls.
With circuit preprocessing, the total costs for the batch are:
\begin{equation}
\begin{aligned}
    \text{Prover time}\quad
       &O(N\kappa^2)\,\bb+\mathcal{TP}_{\mathsf{pcs}}(4N),\\
    \text{Prover memory}\quad
       &O(N\kappa)\,\Field+\mathcal{SP}_{\mathsf{pcs}}(4N),\\
    \text{Verifier time}\quad
       &O\left(\kappa\log_2N+(\log_2N)^2
          +n_{\mathsf{io}}\right)\,\ee+\mathcal{TV}_{\mathsf{pcs}}(4N),\\
    \text{Proof size}\quad
       &O\left(\kappa+(\log_2N)^2\right)\,\ExtensionField+\mathcal{S}_{\mathsf{pcs}}(4N),
\end{aligned}
\label{eq:together-multiple-efficiency}
\end{equation}
where \(\mathcal{TP}_{\mathsf{pcs}}\),
\(\mathcal{SP}_{\mathsf{pcs}}\),
\(\mathcal{TV}_{\mathsf{pcs}}\), and
\(\mathcal{S}_{\mathsf{pcs}}\) denote the aggregate PCS prover
time, prover memory, verifier time, and proof size for total
logical length \(4N\), respectively.

Excluding the underlying PCS, dividing the prover time and peak
memory by \(\kappa\)
gives amortized costs of \(O(N\kappa)\bb\) and
\(O(N)\Field\) per instance.
\end{theorem}

The proof is deferred to
Appendix~\ref{appx:proof-multiple-instance}.
Algorithm~\ref{alg:multiple-instance-packing} gives the protocol.

A PCS whose prover time and memory are linear in its
extension-field input length therefore gives an end-to-end
bitwidth-aware construction.

\iffalse
\subsection{Binary-Field Instantiation}
\label{sec:binary-packing}

For binary circuits, we use a small binary extension as the working
field of the entire PIOP. By contrast, the intermediate-field
Hadamard variant in Section~\ref{sec:hadamard} retains the
original base-field witness and its fixed-basis commitment,
using the intermediate field only for certificate construction.
For compatible \(t\mid\kappa\) and \(2^t\geq2\kappa-1\),
its tabulated implementation has \(O(N\kappa)\) online
base-field cost when \(N\geq\kappa^2\). The storage includes
\(O(\kappa^2t+t2^{2t})\) precomputed base-field coordinates.
The working-field instantiation here supplies interpolation nodes
and represents the witness in byte-field coordinates. Choose:
\[
    \mathbb F_2\subset\mathbb K=\mathbb F_{2^8}
       \subset\ExtensionField,
    \qquad [\mathbb K:\mathbb F_2]=8,
    \qquad [\ExtensionField:\mathbb K]=\kappa_{\mathbb K}.
\]
Arithmetize the circuit over \(\mathbb K\), including constraints
\(w^2-w=0\) for the binary wires, and define \(m,N\) for this
arithmetization. Instantiate the building blocks, packed PCS, and
multiple-instance composition with
\((\Field,q,\kappa)=(\mathbb K,256,\kappa_{\mathbb K})\).
Here \(N=m\kappa_{\mathbb K}\) counts working-field entries.

\smallskip\noindent\textbf{Packing and openings.}
Assume \(256\geq2\kappa_{\mathbb K}-1\), so \(\mathbb K\)
contains the interpolation nodes and off-domain samples needed
by Section~\ref{sec:hadamard}. Represent \(\ExtensionField\) in a
\(\mathbb K\)-basis. All coordinate matrices in
Section~\ref{sec:pcs} use this basis. The verifier checks public
query contractions using Eq.~\eqref{eq:pcs-query-recurrence},
while the gate and copy challenges lie in \(\ExtensionField\).

\smallskip\noindent\textbf{Position encoding.}
Set \(s=8\) and \(\underline q=256\), and assume
\(\lceil\log_{256}L\rceil\leq\kappa_{\mathbb K}\).
Choose an \(\mathbb F_2\)-basis
\((\delta_1,\ldots,\delta_8)\) of \(\mathbb K\) and encode each
eight-bit position block as \(\sum_{t=1}^8\delta_tb_t\).
These digits give injective labels in \(\ExtensionField\) through
Eq.~\eqref{eq:together-multiple-labels}.
The proof in Appendix~\ref{appx:proof-position-encoding} applies
to these coefficients and gives the same MLE identity,
Eq.~\eqref{eq:position-encoding-evaluation}.
Thus the randomized copy argument operates in the stated tower.
The role of \(\mathbb K\) here is to support the PIOP arithmetic.
Using \(\mathbb F_{2^8}\) instead for the joint value-and-position
encoding of Section~\ref{sec:plonk-packing} over \(\mathbb F_2\)
would impose \(1+\lceil\log_2L\rceil\leq8\), or \(L\leq2^7\).

\smallskip\noindent\textbf{Using a binary linear code.}
Appendix~\ref{appx:pcs} gives the code-switching
instantiation of the PCS. A binary generator matrix extends
first to \(\mathbb K\) and then to \(\ExtensionField\) without
changing its entries. The source encoding can therefore be
computed on the eight bit-coordinate matrices over \(\mathbb F_2\),
while the packed opening reduction uses the
\(\ExtensionField/\mathbb K\) basis.
Code switching supplies the terminal extension-field opening
through a characteristic-two-compatible BaseFold
instantiation~\cite{eurocrypt/DiamondP26}.

\smallskip\noindent\textbf{Efficiency and soundness.}
The proof of Theorem~\ref{theorem:multiple-instance-packing} uses
field identities, injective labels, and polynomial degree bounds,
which apply in characteristic two. With the packed PCS of
Section~\ref{sec:pcs}, its efficiency and soundness bounds hold with
\(q=256\) and \(\kappa=\kappa_{\mathbb K}\), subject to the
stated field-size and packing conditions.
Each packed opening has error at most
\(\epsilon_{\mathsf{PCS}}+(1+2\log_2m)/|\ExtensionField|\).
Arithmetic costs are measured in \(\mathbb K\)-operations,
and communication and space in \(\mathbb K\)-coordinates.
For the code-based instantiation, substitute the underlying PCS
error and costs of the chosen Blaze construction as specified in
Appendix~\ref{appx:pcs}. The packed-opening, gate,
Hadamard, and copy-reduction bounds apply with the same parameters.
\fi
